
\فصل{روند جلسات و ارتباطات}

\قسمت{ساختار جلسات}

تیم ساکنا از چارچوب \lr{Lean DAD} پیروی می‌کند و جلسات زیر را در هر ایتریشن برگزار می‌کند:

\زیرقسمت{جلسه‌ی برنامه‌ریزی (\lr{Iteration Planning})}

در ابتدای هر ایتریشن برگزار می‌شود.
\textbf{هدف:} انتخاب آیتم‌های \lr{backlog} برای ایتریشن، تخمین تلاش، تقسیم وظایف بین اعضا.
\textbf{شرکت‌کنندگان:} تمام اعضای تیم.
\textbf{مدت:} ۱ تا ۲ ساعت.

\زیرقسمت{دیلی استندآپ (\lr{Daily Standup})}

هر روز به‌صورت ویس در گروه (تاپیک دیلی) برگزار می‌شود.
\textbf{هدف:} هماهنگی روزانه و رفع موانع.
\textbf{سه سؤال اصلی:}
\شروع{شمارش}
\فقره دیروز چه کردم؟
\فقره امروز چه می‌کنم؟
\فقره آیا مانعی دارم؟
\پایان{شمارش}
\textbf{مدت:} حداکثر ۱۵ دقیقه.

\زیرقسمت{جلسه‌ی دمو (\lr{Demo})}

در پایان ایتریشن اول برگزار می‌شود.
\textbf{هدف:} نمایش نتایج قابل ارائه‌ی ایتریشن به تیم (و در صورت امکان به استاد).
\textbf{شرکت‌کنندگان:} تمام اعضای تیم.

\زیرقسمت{جلسه‌ی رتروسپکتیو (\lr{Retrospective})}

در پایان هر ایتریشن برگزار می‌شود.
\textbf{هدف:} بررسی آنچه خوب پیش رفت، آنچه بهتر می‌توانست باشد و اقدامات بهبود.
\textbf{سه سؤال اصلی:}
\شروع{شمارش}
\فقره چه چیزی خوب بود؟
\فقره چه چیزی می‌توانست بهتر باشد؟
\فقره در ایتریشن بعدی چه کاری انجام می‌دهیم؟
\پایان{شمارش}

\قسمت{کانال‌های ارتباطی}

\شروع{فقرات}
\فقره \textbf{پیام‌رسان:} گروه تیمی در یکی از پیام‌رسان‌های رایج برای هماهنگی روزانه.
\فقره \textbf{مدیریت کد:} \lr{GitHub} برای کنترل نسخه، \lr{Pull Request} و \lr{Code Review}.
\فقره \textbf{مدیریت وظایف:} \lr{GitHub Issues / Projects} برای \lr{backlog} و پیگیری آیتم‌ها.
\فقره \textbf{مستندات:} \lr{GitHub Wiki} یا پوشه‌ی \lr{docs} در ریپازیتوری.
\پایان{فقرات}

\قسمت{جدول زمانی جلسات}

\begin{table}[ht]
\centering
\caption{جدول زمانی جلسات پروژه}
\label{جدول:جلسات}
\begin{tabular}{|>{\raggedleft\arraybackslash}p{3cm}|>{\raggedleft\arraybackslash}p{3cm}|>{\raggedleft\arraybackslash}p{3cm}|>{\raggedleft\arraybackslash}p{3cm}|}
\hline
\rl{\textbf{ایتریشن}} & \rl{\textbf{پلنینگ}} & \rl{\textbf{دمو}} & \rl{\textbf{رترو}} \\
\hline
\rl{ایتریشن ۱} & \rl{خرداد ۲۳} --- \lr{22:00} & \rl{خرداد ۳۱} & \rl{خرداد ۳۱} \\
\hline
\rl{ایتریشن ۲} & \rl{تیر ۱} & \rl{---} & \rl{تیر ۱۲} \\
\hline
\end{tabular}
\end{table}

\قسمت{تعاملات با استاد}

تیم ساکنا برای ارائه‌ی پیشرفت و دریافت بازخورد، در جلسات رسمی درس با دکتر رامان رامسین تعامل دارد.
مستندات تولیدشده (از جمله همین سند چشم‌انداز) در هر مرحله به‌روزرسانی شده و ارائه می‌شوند.
